
% --- Overwrite and append this section to the current Parasite document ---
\section{Appended Refinement: Objection to the Simplex, Prior Literature, and Field-First Revision}

\subsection*{Statement of the objection}

A serious objection arises against the earlier formulation:
\[
C \triangleright (A \to B).
\]

This expression was useful as a first approximation, but it retains a hidden
graph ontology. It suggests:
\begin{itemize}
\item that \(A\) and \(B\) are discrete positions,
\item that the relation \(A \to B\) is something like a line between them,
\item and that \(C\) acts on that line while the terms remain relatively stable.
\end{itemize}

\paragraph{Narrative}
This is the ``thumb tacks and red yarn'' problem.
Even if one says that the parasite acts on the relation rather than on the terms,
one may still be presupposing that the terms are already discretely given.

\subsection*{The objection in relation to prior literature}

This objection is not wholly unprecedented. There are at least three strands in the
literature that already push against reading the Serresian simplex as fixed positions plus a link.

\paragraph{First strand: Serres displaces stations in favor of relations}
Serres himself repeatedly displaces stations in favor of relations.
The key line is the one already isolated:
the parasite has a relation ``with the relation and not with the station.''
He also says that the parasite puts the relation into the form of a cantilever.
This already undercuts any fully substantial reading of \(A\) and \(B\) as self-sufficient points.

\paragraph{Narrative}
So the objection is not invented from outside Serres.
It is already latent in his own emphasis on relation over station.

\paragraph{Second strand: criticism of discrete-entity imagery}
Secondary scholarship also explicitly warns against thinking in terms of discrete entities.
Christopher Watkin's discussion of Serres's ``prepositional thinking'' says that this mode
of thought positions us not as manipulators of ``discreet entities as if they were counters
on a gaming board,'' but as internal interrupters of flows whose actions have distributed effects.

\paragraph{Narrative}
This is very close to the present dissatisfaction with the thumb-tacks-and-yarn image.
It does not yet give a topology-first formalism, but it absolutely raises the same pressure:
fixed nodes are the wrong picture.

\paragraph{Third strand: system scale and blurred node/edge distinction}
Kockelman's treatment of the parasite makes the same issue visible in a more systems-theoretic way.
He distinguishes agents acting on already-existing relations between nodes from agents immanent
to assemblages, and then shows that Serres's parasite stands between these functions,
with what is external becoming internal and vice versa depending on system scale.

\paragraph{Narrative}
This does not fully dissolve nodes into field-effects, but it explicitly problematizes
a clean node/edge distinction and introduces the scale sensitivity on which the present objection depends.

\paragraph{A weaker but relevant strand}
There is also a related line in Serres-adjacent work that treats the parasite as an intermediate
position generating subjectivity only through networks of mediation, rather than as a stable third
point in a static triangle. In that usage, the ``position'' is already inseparable from a network
of channels and flows.

\subsection*{Most honest assessment of the literature}

The most honest answer is therefore twofold:

\begin{enumerate}
\item Yes, the objection has been floated in part:
Serres and some of his readers already resist fixed stations, emphasize relations over terms,
and criticize discrete-entity pictures.

\item No, the full revision stated here does not appear to be standard:
I did not find a source that cleanly says the simplex is only a coarse-grained projection,
that vertices should be treated as metastable folds of a field, or that the parasite acts first
on topology and only derivatively on terms and relations.
\end{enumerate}

\paragraph{Narrative}
So the objection is not ex nihilo.
But the field-first reformulation still appears to be an extension rather than a direct restatement.

\subsection*{Why the objection matters}

If \(A\) and \(B\) are themselves metastable effects of a field, then they cannot be
left untouched by a parasitic modification of the relation.

To alter the relation is also to alter the local field conditions from which
\(A\) and \(B\) emerge.

So the following stronger claim becomes necessary:

\begin{quote}
The parasite cannot act only on the relation abstracted from the terms, because
the terms are themselves conditioned by the same local topology that the relation expresses.
\end{quote}

\subsection*{Possible Serresian response}

One possible Serresian answer would be to multiply relays.

That is:
\begin{itemize}
\item the parasite acts on the relation,
\item the altered relation changes the conditions of \(A\),
\item that change itself forms another parasitic relay,
\item likewise for \(B\).
\end{itemize}

\paragraph{Narrative}
This response is plausible and consistent with Serres' general emphasis on mediation,
chains, and cascades.
But even if accepted, it still implies that the simplex should not be treated as
ontologically primary.

\subsection*{Revision: from graph to topology}

The stronger model therefore shifts from a graph picture to a field picture.

Instead of beginning with:
\[
P = (A,B,C)
\]
as if these were already discretely given terms, we begin with a local field:
\[
\mathcal{F}.
\]

The parasite acts first on the field:
\[
\mathcal{F}' = \Psi_C(\mathcal{F}).
\]

Then terms and relation are read off from the deformed field:
\[
A = \alpha(\mathcal{F}), \qquad B = \beta(\mathcal{F}), \qquad R_{AB} = \rho(\mathcal{F}).
\]

\paragraph{Narrative}
This means:
\begin{itemize}
\item the parasite acts on a topology,
\item the apparent relation is one trace of that deformation,
\item the apparent terms are other traces of the same deformation.
\end{itemize}

\subsection*{Consequences of the revision}

The revision yields three important consequences.

\paragraph{1. Terms are no longer primitive positions}
\(A\) and \(B\) are local stabilizations or folds in a field, not self-sufficient stations.

\paragraph{2. Relations are no longer line-like objects}
A relation is not a string between two points. It is a structured pathway of differential coupling within the field.

\paragraph{3. The parasite necessarily has mediate effects on terms}
Even if its primary operation is at the level of differential coupling, the parasite
reconditions the terms because the terms are effects of that same coupling.

\subsection*{Revised formulation}

The earlier formulation:
\[
C \triangleright (A \to B)
\]
should therefore be retained only as a coarse-grained notation.

The stronger formulation is:
\[
C \triangleright \mathcal{F}
\]
with:
\[
A = \alpha(\mathcal{F}), \qquad B = \beta(\mathcal{F}), \qquad R_{AB} = \rho(\mathcal{F}).
\]

\paragraph{Narrative}
The simplex is still useful, but only as a projection:
\[
P = \Pi(\mathcal{F}).
\]
It is a readable local summary of a deeper field deformation, not the primary ontology.

\subsection*{Refined statement}

A more precise statement is therefore:

\begin{quote}
The parasite acts primarily at the level of differential coupling rather than at the level of substantial terms; but because the terms are themselves local effects of that coupling, the parasite necessarily reconditions them as well.
\end{quote}

\subsection*{Conclusion}

This objection does not destroy the Parasite machine.
It clarifies its proper level.

The simplex remains a useful modeling surface.
But the ontologically stronger claim is field-first:

\begin{itemize}
\item topology before graph,
\item deformation before line,
\item local stabilization before fixed position.
\end{itemize}

\subsection*{References for this refinement}

\begin{itemize}
\item Michel Serres, \emph{The Parasite}, especially the formulation that the parasite has a relation ``with the relation and not with the station.''

\item Christopher Watkin, ``How Serres Thinks,'' in \emph{Parrhesia} 40 (2024), especially the discussion of ``discreet entities as if they were counters on a gaming board'' and prepositional thinking from inside the flow.

\item Paul Kockelman, ``Enemies, Parasites, and Noise: How to Take Up Residence in a System Without Becoming a Term in It,'' especially the discussion of differential distances, immanence, and system scale.

\item Raymond L. M. Lee, ``Affinities of Affectivity: Entanglement, Contagion and Exploitation in Smartphone Society,'' \emph{Media Theory}, for a Serres-adjacent use of the parasite as an intermediate position generated through channels and flows.
\end{itemize}

% --- End append ---
